% !TEX root = ../thesis.tex

%=========================================================
\chapter{Introduction}
\label{ch:introduction}
%=========================================================

    % most broad statement on importance of computers
    Modern science, industry, and everyday life now run on computation. As the scope and scale of computation expand, so does the demand for processing power. Meeting that demand carries a physical cost in the energy required to sustain it \cite{iea_energy_ai_2025}. 
    
    % scaleability and Moore's Law
    For decades, computing advanced to the rhythm of Moore's law. Smaller transistors placed ever more functionality on each chip, while reductions in operating voltage limited the accompanying increase in power consumption \cite{moore_cramming_1965,dennard_design_1974}. This strategy no longer guarantees proportional improvements in performance and energy efficiency because power dissipation, fabrication complexity, material properties, and device physics make further scaling increasingly difficult \cite{markov_limits_2014}.
    
    % quantum computers are the shit
    As conventional scaling approaches its practical limits, quantum mechanics offers more than a constraint. It provides a new basis for computation. A sufficiently powerful quantum computer could solve selected problems beyond the practical reach of classical machines. Quantum simulation could transform how molecules, pharmaceuticals, catalysts, and quantum materials are understood and designed \cite{feynman_simulating_1982,mcardle_quantum_2020}. Quantum algorithms could also break widely used public key cryptosystems, forcing a fundamental change in how digital communication is secured \cite{shor_algorithms_1994}. Quantum computers may not replace classical computers, but their impact on science, industry, and everyday life could be revolutionary.

    % super conducting qpits are the most promising quantum computer design
    Superconducting circuits form the hardware basis of several of the world's most advanced quantum processors and provide a principal route toward scalable quantum computing \cite{devoret_superconducting_2013,kjaergaard_superconducting_2020,arute_quantum_2019}. Their fast operation, lithographic fabrication, and compatibility with integrated microwave control have enabled increasingly complex processors and major advances in quantum error correction \cite{google_quantum_ai_error_2025}.
    
    % Limits of superconducting Qbits
    These processors rely on Josephson junctions, typically formed by two superconducting electrodes separated by an ultrathin insulating barrier. Their nonlinearity enables qubit operation, while their dissipation and microwave response originate from microscopic charge transport. The same scalability that makes superconducting circuits attractive also exposes their microscopic limitations. Nonequilibrium quasiparticles tunneling across a Josephson junction can relax the qubit, while defects in amorphous tunnel barriers absorb energy and produce fluctuations in relaxation time \cite{catelani_relaxation_2011,klimov_fluctuations_2018,bilmes_probing_2022}. As processors grow, understanding and controlling these local loss channels becomes essential to their reproducibility and performance.

    \newpage
    % establish the experimental strategy
    % establish the microwave probe
    To connect the circuit and microscopic descriptions, this thesis compares two complementary junction regimes. Oxide tunnel barriers provide the weak transmission limit, while atomic contacts expose transport through discrete channels of finite and high transmission \cite{scheer_conduction_1997,agrait_quantum_2003}. This controlled progression tests where an effective circuit description remains sufficient and where individual charge transfer processes must be resolved. Microwave irradiation extends this comparison into the driven regime. The spacing and amplitude evolution of the resulting structures reveal transferred charges and overlapping transport processes, directly testing transport descriptions under drive.

    % establish the methods
    In this thesis, I investigate aluminum junctions at millikelvin temperatures under dc bias and microwave irradiation. Oxide tunnel barriers provide the weak transmission reference, while a mechanically controllable break junction forms atomic contacts with different channel compositions. Static spectra establish the superconducting energy scales and channel transmissions. Driven spectra then reveal how quasiparticle, multiple Andreev reflection, and Cooper pair contributions evolve with microwave amplitude and transmission.

    %=========================================================
    \subsubsection*{State of the Art and Open Questions}
    %=========================================================

    % static transport
    Static transport provides a well-established reference for driven measurements. In the tunnel limit, quasiparticle current probes the BCS density of states \cite{bardeen_tunnelling_1961,giaever_study_1961}. As the transmission increases, multiple Andreev reflection produces current within the superconducting gap and the characteristic subharmonic gap structure \cite{averin_ac_1995,cuevas_hamiltonian_1996}. Fitting the complete subgap \textit{I--V} characteristic then allows the discrete channel transmissions of an atomic contact to be extracted \cite{scheer_conduction_1997,agrait_quantum_2003}. Coherent Cooper pair transport provides a complementary macroscopic description. The Josephson relations connect the supercurrent to the superconducting phase, while the RCSJ model embeds this nonlinear element in a resistive and capacitive circuit \cite{josephson_possible_1962,stewart_currentvoltage_1968,mccumber_effect_1968}. A junction cannot generally be separated from its electromagnetic environment, even in static transport. The surrounding capacitance and impedance permit energy exchange during charge transfer and modify the measured \textit{I--V} response. In the tunnel limit, \textit{P(E)} Theory connects this energy exchange to quasiparticle and Cooper pair tunneling \cite{falci_quasiparticle_1991,ingold_charge_1992}.

    % addind microwaves
    Microwave driving adds coherent and incoherent responses to this environment dependent baseline. In the tunnel limit, the Tien-Gordon Model describes photon assisted quasiparticle transport as weighted replicas of the static current \cite{tien_multiphoton_1963}. Coherent Josephson oscillations can instead synchronize with the drive and produce Shapiro steps \cite{shapiro_josephson_1963}. At finite transmission, phase locking produces integer and subharmonic Shapiro resonances, while photon assisted multiple Andreev reflection (PAMAR) generates finite-bias sidebands \cite{cuevas_subharmonic_2002}. Photon assisted incoherent Cooper pair tunneling (PAICPT) produces a distinct zero-bias replica pattern without phase locking \cite{roychowdhury_microwave_2015}. These mechanisms can overlap in the same spectrum, making their separation a central interpretive problem.

    % chauvin et al
    Voltage biased aluminum atomic contacts provide a particularly close precedent. Static MAR yields the channel transmissions, while microwave spectra show coherent integer and fractional Shapiro resonances on a PAMAR background. A low-impedance environment supports phase locking, and microscopic driven MAR theory reproduces representative finite-bias conductance traces using the measured transmissions and calibrated drive without adjustable parameters \cite{chauvin_superconducting_2006}. These measurements establish quantitative PAMAR, but do not address its amplitude resolved separation from incoherent Cooper pair transport.

    % STMs et al.
    Josephson STM experiments identify PAICPT through its replica spacing and squared-Bessel amplitude dependence \cite{roychowdhury_microwave_2015}. Across different STM conductance regimes, quasiparticle tunneling, PAICPT, and driven MAR appear in the same spectra. Direct charge dependent dressing of the static full counting statistics (FCS) components fails to reproduce the irradiated MAR response, while a microscopic calculation performs better but retains substantial deviations \cite{kot_microwave-assisted_2020}. Related measurements of Yu-Shiba-Rusinov states further show that separate electron and hole thresholds can make the visible sideband spacing differ from the net transferred charge \cite{peters_resonant_2020}.

    % Carrad et al.
    Gate tunable semiconductor--superconductor nanowire Josephson junctions reveal charge dependent MAR sidebands in microwave amplitude and frequency maps. A phenomenological comparison dresses the measured static current under the piecewise assumption that one MAR order dominates near each threshold \cite{carrad_photon-assisted_2022}.

    % the Gap
    Taken together, previous studies leave an experimental as well as a modeling gap. They demonstrate individual driven transport mechanisms in selected devices and spectra, but do not systematically map coexisting PAMAR and PAICPT across complete microwave amplitude sweeps and independently characterized contact configurations. Microscopic Floquet treatments capture coherent microwave dynamics but are computationally demanding, while phenomenological constructions are easier to apply but rely on measured line shapes and piecewise assignments \cite{kot_microwave-assisted_2020,carrad_photon-assisted_2022}. What is missing is therefore both a systematic comparison across the few-channel regime and a practical description that predicts the driven response from independently characterized static transport.

    %=========================================================
    \subsubsection*{Contributions}
    %=========================================================

    In this thesis, I systematically measure microwave amplitude sweeps across atomic contacts with substantially different channel compositions. Each sweep is a series of \textit{I--V} curves recorded at increasing microwave amplitude. I display these data as amplitude maps, with bias and drive amplitude as the two coordinates and current or differential response encoded by color. The data reveal coexisting PAMAR and PAICPT throughout the few-channel regime and establish how their trajectories and spectral weights evolve with transmission and drive. To explain these maps, I develop a deterministic threshold decomposition. It partitions the simulated static MAR conductance around its known thresholds and reconstructs an effective current for each charge order after direct FCS dressing fails across the full maps. Charge dependent photon dressing and an independently constrained PAICPT contribution then reproduce the dominant trajectories and spectral weight redistribution without selecting structures from the driven measurement.

    The static environmental analysis and the zero-bias replicas support incoherent Cooper pair transport rather than coherent Shapiro response. Tests from moderate transmissions to a nearly ballistic leading channel show that the combined framework captures the principal few-channel microwave response. The framework provides a practical model for testing charge assignments and separating overlapping mechanisms. Whether a microscopic Floquet treatment would improve the remaining agreement is not established here. In the many-channel regime, the controlled description breaks down and an additional subharmonic photon assisted response remains without a unique microscopic explanation.

    In a complementary tunnel barrier experiment, I observed a microwave induced rectification effect that, to my knowledge, has not previously been reported. The local current asymmetry reaches \qty{25}{\percent}, and this strong bias asymmetric response lies beyond the established bias symmetric Tien-Gordon Model. The effect occurs within a specific operating range set jointly by voltage bias, microwave amplitude, and frequency. Together, these three control parameters tune the magnitude and polarity of the rectification. The microscopic mechanism remains open.

    Chapter~\ref{ch:superconducting-transport} develops the theory from BCS quasiparticle tunneling and Josephson phase dynamics to channel resolved MAR and environmental energy exchange. Chapter~\ref{ch:methods} follows the experiment from sample fabrication and the cryogenic MCBJ setup through filtering and the electrical measurement circuit to acquisition, processing, and fitting.

    Chapter~\ref{chapter:experimental-results} orders the measurements by increasing transmission and model complexity. Section~\ref{sec:tunnelbarrier} examines high-resolution photon assisted quasiparticle transport and frequency dependent rectification in two aluminum tunnel barriers. Section~\ref{sec:atomic-contacts} extracts atomic-contact transmissions, develops the combined PAMAR and PAICPT description, tests it across different channel compositions, and concludes with its breakdown and the subharmonic response at high conductance. Chapter~\ref{ch:conclusion} summarizes the established observations, model limits, and unresolved mechanisms.

    The chapters follow a common experimental strategy. Static transport establishes the superconducting energy scales and channel transmissions. Microwave amplitude sweeps then reveal how quasiparticle, MAR, and Cooper pair contributions evolve under drive. Across contacts with substantially different channel compositions, comparison of the complete response with minimal transport models establishes a common description of coexisting photon assisted MAR and incoherent Cooper pair transport. In a complementary tunnel barrier measurement, the amplitude and frequency resolved response reveals a previously unreported tunable rectification beyond the bias symmetric Tien-Gordon Model.
